% ============================================================
%  K. Kroman, »Et Par yderligere Bemærkninger om »Relativitetsprincipet««
%  Fysisk Tidsskrift, 15. Aargang (1916-17), pp. 192-205.
%
%  DIPLOMATIC TRANSCRIPTION (Danish)
%  Status: COMPLETE. Printed pp. 192-205, all 14 pages, transcribed 2026-09-04
%  from the scan described below. No gaps, no duplicates; braces and guillemets
%  balanced (22/22); compiles with 0 errors, 0 missing characters, 0 overfull
%  and 0 underfull boxes. English translation not begun.
%
%  Kroman's reply to the three articles that answered his 1915 attack -- Holst's
%  »Tidsproblemet«, C. E. Walsøe's »Omkring Videnskabens Grænse« and H. M.
%  Hansen's »Relativitetsprincipet« -- and the close of the debate. At p. 193:
%  »Med Forfatteren af den tredie Afhandling, Hr. Dr. H. M. Hansen, er jeg
%  derimod erkendelsesteoretisk set meget uenig.«  Koch's judgement is that it
%  adds nothing new to the 1915 position. See ../relativitetsprincipet/note.md.
%
%  RIGHTS: clear. Kroman died 26 July 1925; Danish copyright expired at the end
%  of 1995. Free to transcribe and to publish.
%
% ------------------------------------------------------------
%  THE WITNESS
% ------------------------------------------------------------
%  Google Books `Tr4ZAAAAIAAJ` -- Fysisk tidsskrift, VOLUMES 13-15 bound
%  together, the University of California copy, digitised 17 January 2007,
%  886 pages, full view. Kept as texts/kroman/scan-ft-13-15.pdf, untracked;
%  shared with ../relativitetsprincipet and ../holst-tidsproblemet.
%
%  Every one of the 14 pages was inspected as an image, doubtful spots at the
%  native 600 dpi 1-bit rendering.
%
%  KB catalogues only the offprint (14 s., »Særtr. af: Fysisk Tidsskrift«),
%  filed in the Særtrykssamlingen with no call number and marked »U. titelbl.«
%  -- no title page. That is why two attempts to have KB produce this text
%  failed. The offprint's 14 pages match pp. 192-205 exactly.
%
% ------------------------------------------------------------
%  PAGE MAP -- VERIFIED 2026-09-04
% ------------------------------------------------------------
%      printed 192-205  =  PDF + 622   (PDF 814-827), uniform.
%
%  Note that the offset is NOT the same as at the front of the same volume,
%  where it is 602: unpaginated advertisement leaves are bound in (PDF 795 is
%  an ad for H. Struers kemiske Laboratorium). Confirmed on content: PDF 815 =
%  printed 193 carries the running head, PDF 828 = printed 206 is »Anmeldelser.«
%
%  p. 192 IS A SHARED PAGE. Kroman's article begins partway down it, under the
%  running head of the article it follows -- P. O. Pedersen, »Den elektriske
%  Bues Elektronteori«. Only Kroman's text is transcribed; the point where it
%  begins is marked in a comment at the site.
%
% ------------------------------------------------------------
%  CONVENTIONS -- SETTLED AGAINST THE IMAGES
% ------------------------------------------------------------
%  * Diplomatic. 1917 orthography exactly as printed; the pre-1948 aa kept,
%    nouns capitalised, nothing modernised.
%  * QUOTATION MARKS: guillemets »...«, \guillemotleft / \guillemotright,
%    22 pairs, balanced.
%  * EMPHASIS is LETTERSPACING, set \emph{}. 42 runs, and none at all on
%    pp. 192-193 or 200. One run sits INSIDE a compound --
%    Bevægelses\emph{ændrings}aarsag (p. 195) -- and two on p. 205 enclose the
%    italic variables A and B.
%  * FOOTNOTES: marked *), one per page on pp. 193, 198, 200 and 205.
%    p. 198 carries a SECOND note marked **) -- the only such page. \thefootnote
%    is fixed to *) in the preamble and the **) is given its mark inside a group
%    at the site, rather than by resetting the counter, so that hyperref's
%    anchors stay unique.
%  * MATHEMATICS: inline only. No displays, therefore no equation numbers.
%    Variables italic; \dfrac; coordinate triples set (O\,x\,y\,z) as printed.
%    The exponent on p. 200 is a^2, not a^3 -- this face's superscript 2 has a
%    looped top; checked against the 1915 article before setting.
%  * The old Danish id-est sign ɔ: (pp. 195, 202) is built from a reflected c;
%    see the preamble.
%
% ------------------------------------------------------------
%  PRINTER'S DEFECTS -- transcribed AS PRINTED, logged at the site
% ------------------------------------------------------------
%    p. 203  »sikkert!,« -- exclamation mark and comma together
%    p. 204  »udtaler ,at« -- the comma set after the word-space, unmatched
%    p. 200  »2 l a« where p. 201 sets the same product »2la«
%  Footnote marks sit inconsistently about the stop, and are set as printed:
%    BEFORE it on p. 193 and for the **) on p. 198; AFTER it for the *) on
%    p. 198 and on p. 205.
%
% ------------------------------------------------------------
%  HOW IT ENDS
% ------------------------------------------------------------
%  No signature and no dateline. The text simply stops, the p. 205 footnote
%  follows, and below it a centred broken-rule tail-piece, transcribed as a
%  short centred rule with a comment. There are no figures or tables anywhere
%  in the article.
%
% ------------------------------------------------------------
%  BUILD
% ------------------------------------------------------------
%  pdflatex path (libertinus + libertinust1math), so the Makefile's default
%  pattern rule builds it. Sandbox compile-test per TRANSCRIPTION-PLAYBOOK.md §5.
% ============================================================

\documentclass[12pt,a4paper]{article}
\usepackage[utf8]{inputenc}
\usepackage[T1]{fontenc}
\usepackage{libertinus}
\usepackage{libertinust1math}
\usepackage{textalpha}   % Greek in text, typed directly
\usepackage{amsmath}     % this debate has real formulas; see CONVENTIONS

\usepackage{graphicx}    % for \reflectbox, below
% The old Danish id-est sign »ɔ:« (pp. 195, 202). T1 has no OPEN O, so it is
% built from a reflected c, as texts/brochner/det-religioese does.
\DeclareUnicodeCharacter{0254}{\reflectbox{c}}

% FOOTNOTES. The printing marks them *), one per page on pp. 193, 198, 200, 205.
% p. 198 carries a second, marked **); that one is given its mark inside a group
% at the site rather than by resetting the counter, which keeps hyperref's
% anchors unique.
\renewcommand{\thefootnote}{*)}
\usepackage[danish]{babel}
\usepackage[protrusion=true,expansion=true]{microtype}
\usepackage{setspace}
\usepackage[margin=1.2in]{geometry}
\usepackage{fancyhdr}
\setlength{\headheight}{15pt}
\usepackage{hyperref}

\hypersetup{
  pdftitle={Et Par yderligere Bemaerkninger om Relativitetsprincipet},
  pdfauthor={K. Kroman},
  hidelinks
}

\pagestyle{fancy}
\fancyhf{}
\fancyhead[LE,RO]{\thepage}
\fancyhead[RE]{\textit{Et Par yderligere Bemærkninger}}
\fancyhead[LO]{\textit{K. Kroman}}

\setstretch{1.3}

\title{\textbf{\textbf{Et Par yderligere Bemærkninger om}\\[2pt]\guillemotleft Relativitetsprincipet\guillemotright}}
\author{K.\ Kroman}
\date{\textit{Fysisk Tidsskrift} 15 (1916--17), pp.~192--205}

\begin{document}
\maketitle

% --- p. 192 ---
% SHARED PAGE. Printed p. 192 belongs for its upper two thirds to the
% PREVIOUS article, P. O. Pedersen, »Den elektriske Bues Elektronteori«,
% whose running head stands at the head of the page and whose closing
% bibliography (Monash, Child, Rasch, Hagenback, Poulsen, Pedersen) fills
% it down to about mid-page. Kroman's article begins BELOW that, roughly
% three fifths of the way down, with its own centred title block
%     Et Par yderligere Bemærkninger om
%     »Relativitetsprincipet«.
%              Af
%          K. Kroman.
% and the body opens »Fysisk Tidsskrift indleder sin 15. Aargang«.
% Nothing of Pedersen's article is transcribed. There is no rule between
% the two items -- only white space.
% The title block is set by \maketitle in the preamble and is NOT repeated
% here. NOTE: the printing ends the title with a FULL STOP
% (»Relativitetsprincipet«.) and the byline reads »K. Kroman.«; the
% preamble's \title/\author drop both stops.

Fysisk Tidsskrift indleder sin 15. Aargang med ikke mindre end tre
Afhandlinger om det \emph{Einstein}'ske Relativitetsprincip. Da de alle tre
tager Stilling til min Artikkel om samme Emne i 14. Aargang Side 1--30, vilde
jeg gerne have Lov til at fremføre endnu et Par Bemærkninger om Spørgsmaalet.
% NOTE: no letterspacing anywhere on p. 192 or p. 193. »Einstein'ske« here is
% set in ordinary roman, exactly like every other proper name in the article;
% see the emphasis note in the header. It is transcribed WITHOUT \emph.

Med Forfatteren af den første Afhandling, Hr.\ Redaktør Helge Holst, er jeg
vistnok i ét og alt enig. Ogsaa hans Indledningsbemærkning om de Einstein'ske
Formlers praktiske Anvendelighed kan jeg godt underskrive; ja jeg har endog
udtalt det samme øverst S.\ 24 i min Artikkel.
% BACK-REFERENCE: Helge Holst, »Tidsproblemet«, Fysisk Tidsskrift 15,
% the first of the three replies. See ../holst-tidsproblemet.
% BACK-REFERENCE: Kroman's own »Relativitetsprincipet«, Fysisk Tidsskrift 14
% (1915-16), pp. 1-30, here cited as »14. Aargang Side 1-30« and »S. 24«.

Og heller ikke med Forfatteren af den anden Afhandling, Hr.\ Ingeniør
C.\ E.\ Walsøe, er jeg i Grunden, hvad Problemet Einstein angaar, uenig. Naar
han (nederst S.\ 15) finder, at Einsteins Lære er fældet som videnskabelig
Hypotese, er jeg fuldstændig tilfredsstillet. At kalde den en uvidenskabelig
Hypotese har jeg ikke noget imod.
% BACK-REFERENCE: C. E. Walsøe, »Omkring Videnskabens Grænse«, Fysisk
% Tidsskrift 15, cited at its p. 15.

% --- p. 193 ---
Med Forfatteren af den tredie Afhandling, Hr.\ Dr.\ H.\ M.\ Hansen, er jeg
derimod erkendelsesteoretisk set meget uenig. Dr.\ H. optræder som en varm
Forsvarer af Einstein's Lære, som han ikke finder videre forskellig fra
Lorentz's, og han mener, at jeg har misforstaaet denne sidste. Jeg antager, at
dette navnlig skal betyde, at jeg ikke tilstrækkelig har skelnet mellem
Lorentz's oprindelige og hans senere, nærmere udformede Hypotese. Men dette
gav min Artikkel mig slet ikke Anledning til. Dr.\ H. overser, at jeg
udtrykkelig har fremhævet, at jeg ikke vilde skrive en Elektricitetsteori, men
blot undersøge den Einstein'ske Hypotese særlig fra dens erkendelsesteoretiske
Side. Jeg vilde heller ikke give en Fremstilling af de Lorentz'ske Arbejder,
saa meget mindre, som jeg jo ikke er afgjort sikker paa, at de netop giver os
den rette Løsning paa den Michelson'ske Gaade\footnote{Det er virkelig
vedblivende mit Haab, at M.'s Resultat nok en Gang skal blive forklaret ved en
Ændring af Lysteorien. Dr.\ H. finder (S.\ 18), at denne Udsigt er spærret ved
en Ytring af Poincaré. Mig forekommer det imidlertid, at denne Ytring --- i det
mindste i Dr.\ H.'s Formulering --- er ganske uberettiget. Med hvilken Ret
skulde man kunne forlange, at L.'s Hypotese skulde tage Hensyn til andre
Mærkeligheder end M.'s, før der endnu havde vist sig andre? Opdukker der andre
af samme Natur, vil samme Hypotese jo ogsaa kunne forklare dem, og er de af en
ny Natur, maa der vel ogsaa nye Hypoteser til. Men man skal som bekendt ikke
sælge Skindet, før man har skudt Bjørnen.}.
% FOOTNOTE MARK POSITION: the printing sets »Gaade*).« -- the mark BEFORE the
% full stop. Compare pp. 198 and 205, where it stands after it.
% BACK-REFERENCE: H. M. Hansen, »Relativitetsprincipet«, Fysisk Tidsskrift 15,
% cited here at its p. 18 (and at p. 22 in the footnote to p. 198 below).
% See ../hansen-svar.

Lorentz har fremstillet begge sine Hypoteser --- den sidste ikke mindre end
tre Gange --- i de Skrifter, jeg har anført. Han har her omtalt de forskellige
med Michelson's beslægtede Forsøg, som ogsaa Dr.\ H. nævner. Jeg kan altsaa
ikke godt have overset, at han havde baade en oprindelig og en mere udført
Hypotese. Men fælles for dem begge er Forkortningen, og andet og mere behøvede
jeg ikke til min Undersøgelse af Einstein's Tankegang. Jeg har for øvrigt ogsaa
fremhævet Formelfællesskabet mellem Einstein og Lorentz, men samtidig betonet,
at der til Trods herfor dog er en særdeles væsenlig Forskel mellem de to. Jeg
skal i det følgende endnu en Gang søge at begrunde, at Lorentz's Hypoteser i
ethvert Tilfælde dog er mulige Antagelser, medens Einstein's efter min
Opfattelse maa stemples som umulig.

Dr.\ H. kalder Einstein's Lære simplere og (i en tidligere Afhandling)
dristigere end L.'s. At Simpelheden maa være af en ret
% The foot of p. 193 carries the signature line »Fysisk Tidsskrift. XV« with
% the sheet number »14« at the outer edge; p. 195 carries »14*«. Signature
% marks are not text of the work and are not transcribed.

% --- p. 194 ---
overfladisk Art, synes imidlertid tilstrækkeligt at fremgaa af de ualmindelig
kunstlede Resultater angaaende Rum og Tid, den udmunder i, og Dristighed er jo
ikke nogen just ublandet videnskabelig Dyd. Hos Einstein optræder den, som vi
skal se, nærmest som Mod til Overdrivelser og vage Begrebsbestemmelser, Mod til
at trodse al rimelig Sandsynlighed og Ligegyldighed overfor fundamentale
erkendelsesteoretiske Krav. Jeg har efter mit Skøn oplyst dette sidste Punkt
paa adskillige Steder i min Afhandling. Dr.\ H. har søgt at fjerne nogle af de
paapegede Vanskeligheder; men det forekommer mig ikke, at det virkelig er
lykkedes ham. Jeg skal imidlertid ikke dvæle yderligere ved disse Enkeltheder,
da jeg tror at kunne belyse og begrunde Berettigelsen af min Opfattelse endnu
tydeligere ved kritisk at følge Einstein's Tankegang gennem dens Udvikling og
paapege, hvor dens Urigtigheder efterhaanden kommer frem og giver Anledning til
de forskellige Modsigelser og store Usandsynligheder, hvori hans Lære udmunder.

Først maa jeg dog endnu tillade mig et Par almene Bemærkninger.

Dr.\ H. forsvarer forskellige Særegenheder i Einstein's Fremstilling med den
Udtalelse, at Einstein er Fænomenist. Det er imidlertid et lidet betryggende
Forsvar; thi det er alt andet end heldigt for en Forsker at være Fænomenist.
Man forstaar ved en Fænomenist en Forsker, som har dannet sig den Opfattelse,
at det er mest videnskabeligt at holde sig til Fænomenerne, til det,
\guillemotleft som viser sig\guillemotright, og lade Aarsager, Kræfter og
deslige passe sig selv. Det er altsaa Sansningen, der hermed er sat i
Højsædet, mens Tanken som en upaalidelig Raadgiver er sat i Krogen. Men
selvfølgelig er en saadan Erkendelsesteori lige saa umulig at gennemføre, som
den er falsk. Uden Tankens Hjælp vilde der jo hverken kunne regnes eller
sluttes. En virkelig Fornægtelse af Tanken bliver derfor en Umulighed for
Mennesket, og enhver Fænomenist bliver derfor med Nødvendighed en
\emph{inkonsekvent} Forsker, der lader Tanken spille med som sædvanlig de 99
Gange, mens den pludselig den 100. Gang ganske vilkaarligt fornægtes, saa at
man f.\ Eks.\ tillader sig at opstille Virkninger uden at angive eller kunne
angive Spor af Aarsager, og deslige. Af denne Art Skridt vil vi træffe
adskillige hos Einstein, og at slige Friheder ikke er heldige for
Videnskabeligheden, behøver sikkert ikke at paavises.
% FIRST LETTERSPACED RUN OF THE ARTICLE: »i n k o n s e k v e n t«, verified
% at 600 dpi against the native 1-bit scan image. No other run on this page.
% The stray hyphen the text layer reports before »Det er altsaa Sansningen«
% is a speck in the left margin, not an em dash: verified on the image.

% --- p. 195 ---
Lad os da begynde med at se lidt nærmere paa det lidet klare Begreb og det
lidet heldige Udtryk \guillemotleft Relativitetsprincipet\guillemotright, og
lad os først spørge: Havde Newton (og den ældre Fysik i det hele) noget, som
med Rette og Rimelighed kunde kaldes et Relativitetsprincip? Hvad er
overhovedet et Princip?

I de fleste fysiske Lærebøger finder man forskellige ganske skikkelige
Læresætninger haardnakket betitlede Principer: \guillemotleft Arkimedes'
Princip\guillemotright, \guillemotleft Dopplers Princip\guillemotright,
o.\ s.\ v. Dette hidrører aabenbart udelukkende fra sproglig Ligegyldighed og
uheldig Inerti. Ved et Princip forstaas nøjagtigere talt et Udgangspunkt, en
Grundsætning, ɔ: en Sætning, der lægges til Grund for hele den følgende
Lærebygning, en Førstesætning, der ikke er afledet af tidligere, ikke har Bevis
foran sig, idet man har fundet den selvindlysende rigtig, og ingen har vovet at
tvivle om dens Gyldighed. Af saadanne Principer eller Grundsætninger har Newton
som bekendt opstillet tre: Inertisætningen, Kraftsætningen og
Vekselvirkningssætningen. En \guillemotleft
Relativitetsgrundsætning\guillemotright{} har han derimod ikke opstillet. Ved
Kraftsætningen og allerede tidligere har han defineret en Kraft, ikke som en
Bevægelsesaarsag, men som en Bevægelses\emph{ændrings}aarsag, altsaa ikke ved
Hjælp af Bevægelse, men ved Hjælp af Bevægelsesændring, d.\ e.\ Ændring af
Hastighed, Retning eller begge Dele. Den logiske Konsekvens heraf er
bl.\ a.\ den, at har et Legeme en vis Bevægelse, og har et andet lignende den
samme Bevægelse plus en eller anden jævn retlinjet Hastighed, maa
Kraftligningerne for begge selvfølgelig blive ens; thi uændret Bevægelse kræver
ingen Kraft.
% The old Danish id-est sign »ɔ:« occurs here and again on p. 202. It needs
% \DeclareUnicodeCharacter{0254}{\reflectbox{c}} and graphicx in the preamble;
% see the note at the end of this fragment.
% EMPHASIS: the letterspaced run is »æ n d r i n g s« ONLY, standing inside the
% compound word: the printing sets »Bevægelses æ n d r i n g s aarsag«, with
% the spacing of the run doing the separating. Set here as
% Bevægelses\emph{ændrings}aarsag. Verified at 600 dpi.
% The text layer renders the closing guillemet after »Relativitetsgrundsætning«
% as »<<<«; at native resolution it is ONE « (two chevrons), not two marks.

Endvidere hævder han, at et Legeme kan have allehaande Bevægelser, saa vel
relative som absolutte, og at lige saa vel som der er Forskel paa jævn
retlinjet Bevægelse med Hastigheden 5 og jævn retlinjet Bevægelse med
Hastigheden 10, lige saa vel maa \guillemotleft Bevægelse med Hastigheden
0\guillemotright{} være noget for sig selv. Dette er jo ogsaa ethvert besindigt
Menneske nødt til at antage. Ingen kan tænke sig noget ved Udtrykket Bevægelse,
uden at han mere eller mindre tydeligt samtidig maa tænke sig en Hvilen, og
ingen kan tænke sig en Bevægelse som relativ, uden at han samtidig maa tænke
sig en absolut Hvile eller Bevægelse. Disse Forestillinger staar og falder med
hinanden ligesom Forestillingerne Op og Ned, og i enhver rationel Mekanik
tumler man jo fuldstændig sikkert og naturligt med disse Bestemmelser.

% --- p. 196 ---
\emph{Vidt forskelligt herfra} er imidlertid det Spørgsmaal: Kan vi Mennesker
altid bestemme, hvilken Bevægelse et Legeme i et givet reelt Tilfælde har? Med
absolut Nøjagtighed kan vi selvfølgelig aldrig gøre det. Med Tilnærmelse kan vi
derimod ofte bestemme dets relative Bevægelse. I visse Tilfælde vil vi ogsaa
kunne bestemme, at det har den og den absolutte Bevægelse. Men er Bevægelsen
jævn og retlinjet, vil Tilfældet jo --- saa vidt vi endnu véd --- ikke ved
noget sanseligt Kendetegn adskille sig fra en uændret Bevægelse i samme Retning
med en anden eller slet ingen Hastighed, og derfor vil det, som Newton siger,
ofte være \emph{svært} at bestemme den sande Hastighed.

Hvorfor mon Newton nu har brugt et saa vagt Udtryk? Hvorfor har han ikke
rentud sagt, at overfor uændret Bevægelse er Hastighedsbestemmelse slet og ret
umulig? Naturligvis var han ganske klar over, at her var ingen Hjælp at vente
fra Inertiytringer som ved Hastigheds- eller Retningsændringer. Han havde
ganske vist den \emph{Formodning}, at Solsystemets Tyngdepunkt var absolut
hvilende i Rummet. Men det var kun en Formodning. Han har ogsaa udtalt, at der
maaske slet ikke var noget saadant fast Punkt, vi kunde støtte os til. Snarere
har han derfor vistnok tænkt sig, at Materien, som jo saa nylig havde aabenbaret
sig fra en ny Side i den gaadefulde Gravitation, maaske ogsaa en skønne Dag
kunde røbe \emph{visse} Forskelligheder, eftersom den var i Hvile eller i jævn
retlinjet Bevægelse, og at Fysikerne kunde lære at kontrollere disse
Forskelligheder, saa at en Hastighedsberegning derved blev mulig. Er Lorentz's
Hypotese rigtig, har vi jo i \emph{Forkortningen} netop noget i den antydede
Retning.
% EMPHASIS on this page, all verified at 600 dpi: »V i d t f o r s k e l l i g t
% h e r f r a«, »s v æ r t«, »F o r m o d n i n g« (the FIRST occurrence only --
% the second, »kun en Formodning«, is plain), »v i s s e«,
% »F o r k o r t n i n g e n«, »e k s p e r i m e n t e l t« below.

Da senere Bølgeteorien for Lyset var kommet til Magten, og
Interferensfænomenet var blevet nærmere forstaaet, har sikkert adskillige
Fysikere tænkt sig, at saa snart man blot fik tilstrækkeligt teknisk
Herredømme over Interferensforsøgene, vilde man dermed have et Middel til at
skelne mellem absolut Hvile og jævn retlinjet Bevægelse. Netop dette var jo
Michelson's Antagelse.

Vil man udtrykke sig saa nøjagtigt som muligt, tror jeg derfor, man maa
indskrænke sig til om Newton og de ældre Fysikere at sige:

1) De indrømmede, at de foreløbig ikke formaaede \emph{eksperimentelt} at
skelne mellem absolut Hvile og jævn retlinjet Bevægelse.
% The spaced run »e k s p e r i-/m e n t e l t« is broken over a line turn in
% the printing; it is one run and is joined here.

% --- p. 197 ---
2) Men de var ingenlunde sikre paa, at dette vedvarende vilde være umuligt.

3) Endnu mindre opstillede de Umuligheden deraf som et Princip, d.\ e.\ som en
Sætning, man \emph{skulde} gaa ud fra.

4) Og de har sikkert alle været enige med Newton om, at absolut Hvile og jævn
retlinjet Bevægelse var noget i sig selv forskelligt og ikke simpelthen
Afhængigheder af, hvilke Koordinatsystemer man benyttede.

Men forholder det sig saaledes, bliver det ikke blot vildledende, men ganske
urigtigt at benytte Udtryk som \guillemotleft et Newton'sk
Relativitetsprincip\guillemotright. Alt er da langt klarere og skarpere sagt
med de simple Ord, at Newton byggede paa Inertisætningen. ---

Mere betegnende end at kalde Newton Relativist vilde det sikkert være at kalde
ham Absolutist. I Einstein's Lære har vi derimod intet ringere end et Forsøg
paa at relativere ikke blot Hvile og Bevægelse, men endog baade Rum og Tid.

Det er navnlig Michelson's Resultat, der fører ham ind paa disse Bestræbelser.
Michelson gik ud fra, at Lyset uafhængigt af Lyskildens Hvile eller Bevægelse
bestandig i Verdensrummet og i Luften havde den konstante Hast $c$, ligesom
han overhovedet stolede paa den gængse Lysteori. Det samme gjorde Lorentz og
Einstein, og jeg kan derfor forme de følgende Bemærkninger ud fra Urørligheden
af denne Forudsætning.

Michelson fandt nu ved sit Forsøg, at Lyset i én og samme Tid gennemløb de to
Veje $s$ og $s_{1}$.

For menneskelig Fornuft er hertil da ikke andet at bemærke end, at de to Veje
følgelig \emph{maa} have været lige lange. Dette var ogsaa Michelson's første
Tanke, idet han antog, at Jordens Translation havde været modvirket af hele
Solsystemets Bevægelse. Da man fandt denne Udvej spærret, opstillede Lorentz
sin Forkortningshypotese. Ser vi bort fra, at han jo ogsaa kunde have antaget
den kortere Vej forlænget, og fastholder vi Lysteoriens fuldstændige Rigtighed,
var der som sagt ingen anden Udvej, hvis Fænomenet ikke skulde gaa ind under
det rent irrationelle.

Einstein har sikkert kendt denne Lorentz's første Forkortningshypotese, da han
1905 fremkom med sin Lære. Denne hviler, siger han i Begyndelsen af sin
Afhandling, paa to Forudsætninger: Relativitetsprincipet og den dermed
\guillemotleft kun tilsyneladende uforenelige\guillemotright{}

% --- p. 198 ---
Forudsætning om den førnævnte konstante Lyshastighed. Disse to Forudsætninger
er \emph{tilstrækkelige}, tilføjer han. Om Forkortning siger han ikke et Ord.
Derimod erklærer han det overflødigt at antage en
Æter.\footnote{Lorentz, Einstein, Minkowski: Das Relativitätsprincip,
S.\ 27--28.}
% FOOTNOTE MARK POSITION: printed »Æter.*)« -- mark AFTER the full stop, where
% p. 193 sets it before. Both transcribed as printed.

Angaaende \guillemotleft Relativitetsprincipet\guillemotright{} bemærker han,
at han har faaet den Formodning, at der til Begrebet absolut Hvile hverken paa
mekanisk eller paa elektrodynamisk Omraade svarer nogen Ejendommelighed ved
Fænomenerne. Denne Formodning vil han nu \guillemotleft ophøje til
Princip\guillemotright{} og kalde Relativitetsprincipet.

Men saaledes gaar det virkelig ikke an uden videre at ophøje Formodninger til
\guillemotleft Principer\guillemotright. Foreløbig ved vi kun, at i nogle faa
Tilfælde har der vist sig Resultater svarende til Michelson's. Einstein kan
derfor kun have Lov til at opstille som \emph{Hypotese}, at hans Formodning
gælder alment, og han kan derpaa stille sig som Opgave at prøve denne Hypoteses
Mulighed og eventuelle videnskabelige Værd. Men han faar derved helt andre
Forpligtelser, end hvis han straks kunde ophøje sin Formodning til Princip.

Hans første videnskabelige Forpligtelse vilde saaledes blive den atter at
opgive Hypotesen paa Grund af dens Umulighed. Den er nemlig umulig, idet det jo
staar fast, at naar Lysteoriens Rigtighed forudsættes, er Forkortningen eneste
tænkelige Udvej. Men forkortes Legemerne ved Bevægelsen, saa svarer der jo en
bestemt Ejendommelighed, nemlig Uforkortethed, til den virkelige Hvile. At vi
foreløbig ikke \emph{direkte} kan maale Forkortningen, er praktisk set
kedeligt nok. Videnskabelig set er det imidlertid allerede tilstrækkeligt, at
vi kan godtgøre, at den maa finde Sted, ja endog beregne dens Størrelse i
Forhold til enhver forudsat Hastighed.

Men dermed bliver, som vi skal se, al egenlig Relativering af Hvile og
Bevægelse, Rum og Tid, tilintetgjort og de mange tilsvarende overdrevne Udtryk
forkastelige{\renewcommand{\thefootnote}{**)}\footnote{Ogsaa Dr.\ H. benytter
disse vildledende Udtryk. Gentagne Gange hedder det: Hvile og jævn retlinjet
Bevægelse er et og det samme eller fysisk et og det samme; der er ingen Mening
i at skelne mellem dem. Og saa skelner han dog saa kraftigt imellem dem, at han
hævder, at Bevægelsen faar Legemerne til at trække sig sammen, Ure til at gaa
langsommere, Folk til at ældes langsommere; ja han tilføjer endog (S.\ 22), at
enhver Iagttager kan konstatere dette.}}. Allerede her ses da saaledes, at
Einstein's Hypotese er en Umulighed. Hvad han med videnskabelig
% THE ONLY PAGE IN THE ARTICLE WITH TWO FOOTNOTES. The printing marks them
% *) and **). \thefootnote is fixed at *) in the preamble, so the second is
% given its mark inside a group here rather than by resetting the counter --
% that keeps hyperref's anchors unique. Printed »forkastelige**).« -- mark
% BEFORE the full stop, the opposite of »Æter.*)« four paragraphs above.
% The footnote texts carry no letterspacing and no guillemets, though the
% body would take them: verified at 600 dpi.

% --- p. 199 ---
Ret kunde have opstillet og søgt at begrunde, er noget vidt forskelligt fra,
hvad han har udtalt. Han kunde med Rette have opstillet den Hypotese, at der er
en vis \emph{teknisk} Ækvivalens mellem absolut Hvile og jævn retlinjet
Bevægelse.

Uden at mærke noget til Modsigelsen gaar Einstein imidlertid videre. Han
tænker sig to Koordinatsystemer med indbyrdes parallelle Akser og fælles
$X$-akse, det ene $(O\,x\,y\,z)$ hvilende, det andet $(O'\,x'\,y'\,z')$ med
den konstante $X$-hastighed $v$. I hvert er en Iagttager ($A$ og $B$) forsynet
med fuldkomne ens Ure og Maalestokke og fuldendt Evne til at bruge dem. De
stoler begge paa Lysteorien og den konstante Lyshastighed $c$. Dr.\ H. kalder
dem Fysikere og antager, at de kender Michelson's Forsøg. Alt dette kan vi uden
Betænkning forudsætte.
% MATHEMATICS: the printing sets the coordinate letters in italic with a thin
% space between them and the parentheses in roman -- (O x y z), (O' x' y' z').
% Reproduced with \, ; verified at 600 dpi.

Men Einstein vedbliver saa: De antager begge, at de hviler.

Dette synes at være en mærkelig letsindig og dogmatisk Antagelse at tildele
dem. Fysikere, der er saa fortrinligt udrustede, burde vel snarere kritisk have
gentaget Michelson's Eksperiment og i ethvert Tilfælde have bygget paa den
rimeligere Forudsætning, at enhver af dem havde en jævn retlinjet Bevægelse med
en vis ubekendt Hastighed (maaske $= 0$ for den ene af dem), de enten maatte
bevare som ubekendt i deres Regning eller se at faa bestemt saa godt, det lod
sig gøre, f.\ Eks.\ ved Hjælp af en Fiksstjernehimmel eller, hvad de nu havde
som bedst brugbart.

Vi kan imidlertid tænke os, at Einstein har tildelt dem disse vilkaarlige
Antagelser i bestemt Øjemed, maaske for ret tydeligt at vise, at Hvile og jævn
retlinjet Bevægelse er ganske ens eller deslige, og naturligvis har han Lov til
i sit Tankeeksperiment at indføre hvilke som helst Bestemmelser, blot han
meddeler Læseren, hvad han gør.

Men samtidig har Læseren da naturligvis ogsaa Lov til at mærke sig, hvad der
sker, Lov til at mærke sig, at her indførte Einstein en Modsigelse i
Situationen, idet han gav $A$ en rigtig Forudsætning og $B$ en falsk, i direkte
Modsigelse med, hvad der blev sagt om de to Koordinatsystemer.

$A$ ordner derpaa sine rundt i Systemet $O$ opstillede Ure ved Hjælp af
Lyssignaler, saa at de alle faar samme Stand og rigtig Gang, 1 Sekund, mens
Lyset løber Vejen $c$. Dette gaar nu let nok. Men saa skal $B$ gøre det samme
ud fra sin modsigelsesfulde Forudsætning.

% --- p. 200 ---
Og her maa man da huske paa, at Einstein hidtil ikke har sagt et Ord om
Forkortning, hvad han jo heller ikke godt kunde, da han tvertimod har sagt, at
Fænomenerne skulde forholde sig ens under Hvile og jævn retlinjet Bevægelse. Vi
maa derfor foreløbig regne uden nogen Forkortning. Einstein har jo mulig ment,
at den vil indfinde sig af sig selv som Konsekvens af hans Forudsætninger.

Men sæt nu først, at de \guillemotleft fuldkomne\guillemotright{} Ure ogsaa har
haft rigtig Gang (1 Sekund, mens Lyset virkelig løber Vejen $c$) allerede fra
Fabrikantens Haand og at $B$ har stolet herpaa saa vel som paa den konstante
Lyshast. Hvad vil saa ske?

$B$ vilde da simpelthen opdage, at sendte han et Lyssignal Vejen $l$ frem og
tilbage vinkelret paa $X'$-aksen, vilde Lyset ikke være $\dfrac{2l}{c}$
Sekunder, men $\dfrac{2l}{c}\,a$ Sekunder\footnote{Sml.\ angaaende dette og de
følgende Udtryk min tidligere Artikkel.} om Vejen, og sendte han Signalet Vejen
$l$ frem og tilbage langs $X'$-aksen, vilde det bruge Tiden
$\dfrac{2l}{c}\,a^{2}$. Men deraf vilde $B$ sikkert slutte: Mit System kan
altsaa ikke hvile; det maa have den jævne retlinjede Hast $v$.
% MATHEMATICS. Every fraction on pp. 200-201 is a built-up two-line fraction
% standing in the run of the text, so \dfrac is used, as in the 1915 article.
% The exponent in »a^2« was checked at native (600 dpi, 1-bit) resolution and
% against the same journal's superscript 2 in the 1915 article: this face's
% superscript 2 has a closed, looped top and reads as a 3 at lower resolution.
% It is a 2, and a^2 is also what the physics requires (a = the Lorentz factor,
% cf. »Relativitetsprincipet« p. 10). NOT a printer's error.
% There are NO displayed and NO numbered equations anywhere in this article.

Antager vi dernæst, at $B$ ikke stoler paa sine Ures rette Gang, men selv vil
indstille dem, faar vi et lidt ændret Resultat. Lad ham igen begynde vinkelret
paa $X'$-aksen og atter bruge den afmaalte Vej $l$. Indretter han nu sit Ur til
at gaa $\dfrac{2l}{c}$ Sekundstreger frem, mens Lyset er undervejs, vil hans
Sekunder blive $a$ Gange for store; thi Lyset har da i Virkeligheden løbet
Vejen $2\,l\,a$. Hans Ur kommer altsaa til at gaa $a$ Gange for langsomt, hvad
han dog foreløbig ikke aner. Men sender han derpaa Signal langs $X'$-aksen, vil
han til sin Overraskelse finde, at det regulerede Ur nu er gaaet
$\dfrac{2l}{c}\,a$ \guillemotleft Sekunder\guillemotright{} frem, mens Lyset
var undervejs. Og har han tænkt sig lidt om, vil han sikkert ogsaa heraf
slutte, at han ikke blot har den førnævnte Hast $v$, men at han endvidere har
indrettet sit Ur til at gaa $a$ Gange for langsomt.
% The printing sets this product with visible spaces, »2 l a«, where p. 201
% sets the same quantity closed up as »2la«. Both reproduced as printed.

Saaledes maa Resultaterne forme sig, naar vi regner uden Forkortning.

% --- p. 201 ---
Men disse Resultater kan Einstein ikke bruge; thi de staar i Strid med det
Michelson'ske Resultat.

\emph{Kunde} han indføre en real Forkortning, saa kunde $B$'s Selvbedrag
derimod lykkes. Blev alle $B$'s Afstande i $X'$-retning $a$ Gange saa smaa,
vilde han kunne indrette alle sine Ure til at gaa $\dfrac{2l}{c}$
Sekundstreger frem, mens Lyset ud fra hans Forudsætning løb Vejen $2l$, men i
Virkeligheden Vejen $2la$. Hans Ure kom derved til at gaa $a$ Gange for
langsomt, og hans urigtige Forudsætning vilde endvidere medføre, at de fik
Stedtid eller i Forhold til Uret i $O'$ Klokkesletsunderskud proportionale med
deres $X'$-afstande fra $O'$, uden at han vilde mærke noget. Alt dette faas let
ved ganske elementær Regning, og paa samme Maade faas let, hvad $A$ og $B$ i
det hele vilde føres til at mene om hinandens Beregninger, idet de fastholdt
deres indbyrdes stridende Forudsætninger.

Saaledes vilde vi altsaa kunne komme til \guillemotleft de Einstein'ske
Transformationsformler\guillemotright.

Men Einstein selv kan ikke (retmæssigt) komme til dem. Thi han kan hverken
\emph{taale} eller \emph{undvære} eller \emph{indføre} Forkortningen.

Han kan ikke \emph{taale} Forkortningen; thi den vilde tilintetgøre hans
første Forudsætning, der staar i Strid med den, idet den ikke tillader Forskel
mellem Fænomenerne under absolut Hvile og under jævn retlinjet Bevægelse.

Men selv om Einstein ogsaa kunde \emph{taale} Forkortningen, kan han dog ikke
\emph{indføre} den. Thi hvor skulde han faa den fra? Da han har forkastet
Æteren, vilde han slet ikke kunne finde nogen Aarsag til denne Forandring af
Fænomenerne. Forkortningen vilde blive et rent mirakuløst Fænomen, og man har i
Videnskaben dog vel ikke Lov til at opdigte Mirakler for at faa afstivet en
ellers faldefærdig eller uigennemførlig Hypotese?

Paa den anden Side har vi nylig set, at Læren slet ikke kan \emph{undvære}
Forkortningen. Den er ganske nødvendig, for at $B$ i Længden skal kunne
forblive i sin falske Antagelse. Og sker dette ikke, vil Einstein ikke blot
komme i Strid med Michelson; men hans Tankegang vil ligefrem komme til at
godtgøre det modsatte af, hvad han gik ud paa at bevise.

% --- p. 202 ---
Der er derfor næppe andet at gøre end at opgive den Einstein'ske Lære som
\emph{umulig}.

Maaske vil en og anden af hans Tilhængere bemærke, at er der ikke mere i
Vejen, saa staar Sagen ikke saa daarligt for Einstein. Thi vi behøver altsaa
blot at indskrænke hans Relativitetsforudsætning lidt, lade ham beholde Æteren
og lade ham tilføje Forkortningen; saa er alt i Orden.

Ja, ganske vist. Men saa er det blot ikke længer Einstein's Lære; saa er det
simpelthen Lorentz's, vi er kommet til. Og Lorentz's Opfattelse er som sagt en
Mulighed, men til Gengæld, som vi skal se, rigtignok meget forskellig fra
Einstein's.

Vanskeligheden med Forkortningen har Einstein aabenbart slet ikke klart haft
Øje for. Idet han skal til at beregne $B$'s \emph{bevægede} System, gaar det
ganske let for ham, \guillemotleft indem man das Prinzip der Konstanz der
Lichtgeschwindigkeit im ruhenden Systeme anwendet.\guillemotright{} Men det har
han naturligvis slet ikke Lov til. Han begaar derved en Cirkelslutning og faar
slet intet bevist eller højst, at Hvile og Bevægelse er et og det samme, naar
man blot dristigt forudsætter, at de er et og det samme. Fejlen hænger sammen
med, at han uden videre \guillemotleft ophøjer\guillemotright{} det til
Princip, som kun er Hypotese. Har Lyset ifølge hans anden Forudsætning konstant
Hast i Rummet og altsaa i det hvilende System, saa er det \emph{umiddelbart set
utænkeligt}, at det ogsaa skulde have konstant Hast, ɔ: samme Hast i alle
Retninger, i et bevæget System. Noget saadant bliver først tænkeligt, saafremt
det bevægede System og alt dets Indhold paa Grund af Bevægelsen har trukket sig
sammen efter visse bestemte Love. En saadan Forandring kan Einstein imidlertid
ikke faa indført, uden at han maa opgive sin første Forudsætning og atter
antage Æteren. Hans to Forudsætninger staar altsaa \emph{virkelig} i
Modsigelse med hinanden og ikke blot tilsyneladende, som han mente.
% The German quotation is set upright, inside guillemets like every other
% quotation in the article -- NOT in italic. Verified on the image.

Einstein kunde med al videnskabelig Berettigelse have indledet sit Problem med
følgende Spørgsmaal:

Kan man uden at komme i Modsigelse med sund Fornuft antage: 1) at Fænomenerne
altid er ens under absolut Hvile og jævn retlinjet Bevægelse, og 2) at Lyset
altid i Rummet (og Luften) har samme konstante Hast $c$?

Og derpaa maatte da svares: Nej! Thi skal Fænomenerne \emph{endog blot
tilsyneladende} være ens under absolut Hvile

% --- p. 203 ---
og jævn retlinjet Bevægelse, saa maa Lyset ogsaa i et jævnt retlinjet
\emph{bevæget} System tilsyneladende have Hasten $c$ i alle Retninger. Men
dette er utænkeligt, saafremt Systemet og alt dets Indhold ikke under
Bevægelsen er forkortet efter visse (f.\ Eks.\ de Lorentz'ske) Love. Og er
dette Tilfældet, er Fænomenerne jo ikke ens under absolut Hvile og jævn
retlinjet Bevægelse. Vilde man endnu indvende, at Forkortningen ikke kan regnes
til \guillemotleft det fænomenale\guillemotright, til det, der \guillemotleft
viser sig\guillemotright, saa glemmer man, at den jo netop har vist sig gennem
Michelson's Interferensstribers Stilling.

Men Einstein ser intet af alt dette. Han beholder derfor sine Forudsætninger
uden at ane noget om deres Modsigelse. Han ophøjer sin umulige Hypotese til
Princip og tager Forkortningen med paa fri Haand. Kan han antage Lysteorien
uden at antage en Æter, vilde det jo ogsaa være urimeligt at gøre Ophævelser
over saadan en Smule Forkortning; er man Fænomenist, volder sligt ingen
Standsninger. Og saaledes kommer han til det Resultat, at han virkelig har
godtgjort, at Hvile og jævn retlinjet Bevægelse baade paa mekanisk og paa
elektrodynamisk Omraade er et og det samme. Thi han har jo \guillemotleft
bevist\guillemotright, at $A$ og $B$ begge kommer til de samme Resultater
angaaende Fænomenerne, ja angaaende hinanden indbyrdes, skønt $A$ gaar ud fra
en efter gammel Tro rigtig Forudsætning og $B$ fra en efter gammel Tro
fuldstændig gal. Følgelig maa det dog være ganske ligegyldigt, om man
forudsætter sig hvilende eller jævnt retlinjet bevæget, og man kan altsaa
simpelthen sige til Michelson: Du skal blot forudsætte, at Jorden staar stille,
saa bliver dit Resultat straks forklarligt. Siger nogen: Ja, men Jorden bevæger
sig dog ganske sikkert!, svarer Du blot: Hvile og jævn retlinjet Bevægelse er
ganske det samme. Der er ingen Mening i at skelne mellem dem.
% PRINTER'S USAGE, transcribed as printed: »ganske sikkert!, svarer Du blot«
% -- an exclamation mark followed by a comma. Verified at 600 dpi.
% Note also that this whole reported dialogue is set WITHOUT quotation marks,
% where the same author quotes Einstein in guillemets elsewhere.

Dette er egenlig Einstein's Forklaring af Michelson.

Men er Hvile og jævn retlinjet Bevægelse saaledes et og det samme, er der jo
heller ingen Mening i at kalde $A$'s Forudsætning rigtig og $B$'s falsk. Saa er
de begge rigtige og Modsigelsen mellem dem noget, der sikkert vil forsvinde ved
en rigtigere Opfattelse af Rummets og Tidens Natur. Og det er altsaa heller
ikke rigtigt, at $B$'s Sekunder er falske. De er lige saa rigtige som $A$'s,
skønt de er $a$ Gange saa store. Sandheden er blot den, at der er mange Slags
Sekunder i Verden, skønt man godt samtidig kan sige, at et Sekund er den Tid,
hvori Lyset med sin konstante Hast $c$ løber Vejen $c$.

% --- p. 204 ---
Dette er heller ingen virkelig Modsigelse; men Rum og Tid er hidtil bleven
opfattede saa urigtigt, at det ser saaledes ud. Saasnart de opfattes i deres
rette Forhold til hinanden, kommer alt i den skønneste Orden. ---

Saaledes udmunder Einstein's Lære. Modsigelserne i den skubbes længer og
længer ud i det fjerne, ligesom man engang sendte Forbrydere til Botany Bay.
Til sidst gaar det ud over selve Grundbestemmelserne Rum og Tid. Dermed kan man
ikke komme længer, med mindre man endnu vilde gøre et allersidste Skridt og
paastaa, at Modsigelse og Ikke-Modsigelse i Grunden ogsaa er et og det samme.

Lad os dernæst se lidt paa Lorentz's Lære.

Man kan som antydet godt komme til denne Lære ved at gaa ud fra Einstein's og
efterhaanden fjerne alle dens Modsigelser og Overdrivelser. Lorentz er
imidlertid gaaet en anden Vej. Idet han \emph{stolede paa Lysteorien} og fandt
det fornødent at antage en hvilende Æter, forklarer han sig Michelson's
Resultat paa den før omtalte eneste endnu mulige Maade: ved at formode, at
Bevægelsen har medført en Forkortning af Apparatet i Translationens Retning. Og
saa snart de beslægtede, af Dr.\ H. nævnte Forsøg var traadt til, søgte han
yderligere at forklare sig Forkortningens Oprindelse paa en saadan Maade, at
netop de vundne Resultater maatte fremkomme. Derved opstod hans mere
udarbejdede Hypotese, hvis Spirer dog allerede antydes i den første. Han
fremsætter sin Antagelse med stor Forbeholdenhed, kalder den bestandig blot en
Hypotese og udtaler ,at man endnu slet ikke kan vide, hvor almindeligt den
vedkommende Forkortning vil indtræde; derom vil kun Erfaringer efterhaanden
kunne belære os. Hans Undersøgelse har blot paavist, at en saadan Forkortning
uden at give Modsigelse \emph{vilde kunne} fremkomme ad de angivne Veje; men
den har ikke paavist, at den \emph{maa} fremkomme. Lorentz opstiller heller
intet Dogme om Lyshasten som den største i Verden; han siger ganske simpelt, at
hans Hypotese kun gælder for Hastigheder mindre end Lysets. Han gaar med andre
Ord bestandig \emph{korrekt videnskabeligt} til Værks.
% PRINTER'S DEFECT, transcribed as printed: »udtaler ,at man endnu slet ikke
% kan vide« -- the comma is set AFTER the word space, tight against »at«,
% where the sense wants »udtaler, at«. There is no matching mark anywhere in
% the sentence (it closes plainly at »belære os.«), so it is a misplaced comma
% and not an opening quotation mark. Verified at 600 dpi on the native image.

Hvad bliver nu Følgerne af hans Antagelser?

Først og fremmest dette, at alle Vegne, hvor den formodede Forkortning af
Tingene paa Grund af deres Bevægelse virkelig træder

% --- p. 205 ---
til --- og det er jo \emph{muligt}, at den \emph{altid} vil træde til ---, der
vil det gaa os, som det gik Michelson: arbejder vi paa elektrodynamisk eller
optisk Omraade i et jævnt retlinjet bevæget System uden at tage Hensyn til, at
det muligt bevæger sig, vil vi paa Grund af Forkortningen tilsyneladende faa
samme Resultater, som vi vilde have faaet, hvis Systemet havde
hvilet.\footnote{Overfor Spørgsmaal, hvor Tiden ogsaa spiller med, vil dette
dog strengt taget kun indtræffe, saafremt vi enten selv indstiller vore Ure til
tilsyneladende rigtig Gang ved Hjælp af Lyssignaler, eller saafremt selve
Forkortningen virkelig vil bevirke, at de faar denne Gang, selv om Systemet
tidligere har hvilet og deres Gang da svarede til Hvile. Dette Spørgsmaal er
endnu næppe undersøgt tilstrækkeligt.}

Endvidere vil Transformationsformlerne mellem to saadanne Systemer som rent
matematiske Konsekvenser naturligvis faa samme \emph{Form} hos Lorentz som hos
Einstein.

Men de vil faa vidt forskellig \emph{Betydning}.

Lorentz er ikke Relativist. Han siger bestandig \guillemotleft det
Einstein'ske Relativitetsprincip\guillemotright, naar han taler om det. Selv
antager han absolut Hvile og uændret Bevægelse for to forskellige Tilstande, og
det i Forhold til Æteren hvilende System er for ham det virkelig hvilende, det
\guillemotleft udmærkede\guillemotright{} System, som Relativisterne fornægter,
fordi vi foreløbig ikke \emph{eksperimentelt} kan paapege det. Er $A$'s System
et saadant, saa er for Lorentz $A$'s Tid den sande Tid og $B$'s kun en
apparent eller betinget, og det tilsvarende gælder naturligvis deres
Maalinger. Der bliver derfor hos Lorentz ikke Tale om den \emph{Modsigelse},
at \emph{baade $A$ og $B$ har Ret}. \emph{Sandheden er hos ham kun én}, selv om
vi ikke altid kan udfinde, hvor den ligger. Vi véd i ethvert Tilfælde, at højst
kun den ene af de to kan have Ret. Lorentz fornægter saaledes hverken Tankens
Krav eller dens Resultater, og han skelner bestandig klart mellem den rationelle
Sandhed og det eksperimentelt opnaaelige. Rum, Tid og Samtidighed bevarer
derfor hos ham uforstyrret deres tidligere Karakter, og hvad der hos Einstein
optraadte som det kunstlede og paradoksale, netop fordi det blev opfattet som
noget absolut eller enegyldigt, fremtræder hos Lorentz som det ganske
forstaaelige og naturlige, idet det tydeligt stemples som Udtryk for den
\emph{sammensatte} Situation, der fremkommer dels ved den medspillende
Forkortning og dels ved den bevidste Bortseen fra Systemets mulige eller
virkelige Bevægelse.
% EMPHASIS on p. 205, all verified at 600 dpi: »m u l i g t«, »a l t i d«,
% »F o r m«, »B e t y d n i n g«, »e k s p e r i m e n t e l t«,
% »M o d-/s i g e l s e« (broken over a line turn), »b a a d e A o g B
% h a r R e t«, »S a n d h e d e n e r h o s h a m k u n é n«,
% »s a m m e n s a t t e«. In the two longest runs the italic capitals A and B
% stand INSIDE the spaced run; they are set as mathematics within \emph.
% Kroman is here answering the sentence of his own 1915 article that Høffding
% later cites, »Sandheden er ikke dobbelt. Sandheden er kun én«
% (Relativitetsprincipet, p. 15).

\begin{center}
\rule{0.16\textwidth}{0.5pt}
\end{center}
% END OF THE ARTICLE. There is no signature, no dateline, no »Fortsættes« and
% no closing rule of the usual sort: the text simply stops after the footnote,
% and a short centred tail-piece follows about two lines below it -- a broken
% rule of graduated dashes, most of it barely inked in this copy, rendered
% here as a plain short rule, exactly as at the end of the 1915 article.
% PDF 828 = printed 206 begins a different item, »Anmeldelser.«
%
% ------------------------------------------------------------
%  WHAT THE PREAMBLE STILL NEEDS  (this fragment does not compile without it)
% ------------------------------------------------------------
%  1. \renewcommand{\thefootnote}{*)}   -- MISSING. The printing marks
%     footnotes *) (pp. 193, 198, 200, 205) and, once, **) (the second on
%     p. 198, which this fragment marks locally inside a group).
%  2. \usepackage{graphicx} and
%     \DeclareUnicodeCharacter{0254}{\reflectbox{c}}   -- MISSING. The old
%     Danish id-est sign »ɔ:« occurs on pp. 195 and 202; T1 has no open o, so
%     it is built from a reflected c, as ../relativitetsprincipet does.
%  Nothing else is needed: amsmath is already loaded, and there are no
%  displays, no equation numbers, no figures and no tables in this article.

\end{document}
